\chapter{Case Studies}

\section*{Patient focus}
This chapter walks through real-world examples of how imaging is used to diagnose common conditions. Each case tells the story of a patient's journey---from symptoms to scan to diagnosis. These examples are for learning purposes and show how doctors think through a problem step by step. If you are a patient, these cases may help you understand why your doctor ordered a particular test and what the results might mean.

\section*{Imaging focus}
Each case is structured with clinical presentation, imaging findings, differential diagnosis, final diagnosis, and a teaching point. Work through each case systematically before reading the answer. Use these examples to develop pattern recognition, differential diagnosis skills, and an understanding of how imaging guides management.
\label{ch:ch98_cases}

\section{How to Use These Cases}
For each case, follow this approach:
\begin{enumerate}
    \item Read the clinical presentation and form a differential diagnosis
    \item Consider which imaging modality is most appropriate
    \item Review the described findings and interpret them
    \item Compare your interpretation with the final diagnosis
    \item Reflect on the teaching point and how it applies to practice
\end{enumerate}

\clearpage

% ========================================================================
\section{Case 1: Acute Appendicitis on CT}
% ========================================================================

\subsection{Clinical Presentation}
A 28-year-old woman presents to the emergency department with 18 hours of progressive right iliac fossa pain, nausea, and low-grade fever (38.2$^\circ$C). Pain started around the umbilicus and migrated to the right lower quadrant. WCC is 14.2 $\times 10^9$/L with neutrophilia. Urine pregnancy test is negative.

\subsection{Investigation}
CT abdomen and pelvis with IV contrast is requested.

\subsection{Imaging Findings}
\begin{itemize}
    \item A 12 mm diameter, non-compressible, tubular structure in the right iliac fossa arising from the caecal base
    \item Wall thickening with periappendiceal fat stranding
    \item A 5 mm appendicolith within the lumen
    \item No free fluid or extraluminal air to suggest perforation
    \item Surrounding mesenteric lymph nodes are prominent
\end{itemize}

\subsection{Differential Diagnosis}
\begin{table}[h]
\centering
\caption{Differential diagnosis for right iliac fossa pain in a young woman}
\begin{tabular}{|p{0.35\textwidth}|p{0.53\textwidth}|}
\hline
\textbf{Diagnosis} & \textbf{Distinguishing features} \\
\hline
Acute appendicitis & Tubular structure $>6$ mm, fat stranding, appendicolith \\
\hline
Crohn's disease & Bowel wall thickening, skip lesions, chronic history \\
\hline
Mesenteric adenitis & Multiple enlarged nodes, normal appendix \\
\hline
Ovarian pathology & Adnexal mass, free fluid, normal appendix \\
\hline
Right ureteric calculus & Hydronephrosis, ureteric stone on non-contrast \\
\hline
\end{tabular}
\end{table}

\subsection{Final Diagnosis}
Acute uncomplicated appendicitis.

\subsection{Management}
Laparoscopic appendectomy performed within 6 hours. Post-operative recovery uneventful.

\subsection{Teaching Point}
CT is highly accurate for acute appendicitis in adults, although performance varies by protocol and population. Diagnosis uses the overall pattern, including appendiceal enlargement, wall changes and periappendiceal inflammation; a diameter above 6 mm or an appendicolith alone is not diagnostic.

\clearpage

% ========================================================================
\section{Case 2: Pneumonia vs Lung Cancer on CXR}
% ========================================================================

\subsection{Clinical Presentation}
A 62-year-old man, 40 pack-year smoker, presents with a 3-week history of cough, weight loss (4 kg), and intermittent haemoptysis. He has been treated with two courses of antibiotics by his GP with no improvement.

\subsection{Investigation}
PA chest X-ray, followed by CT thorax with contrast.

\subsection{Imaging Findings}
\textbf{Chest X-ray:}
\begin{itemize}
    \item Irregular soft tissue opacity at the right hilum
    \item No cavitation visible
    \item Right hilum appears enlarged and convex
    \item No pleural effusion
\end{itemize}

\textbf{CT thorax:}
\begin{itemize}
    \item 4.5 cm spiculated mass in the right upper lobe abutting the hilum
    \item Associated right hilar and mediastinal lymphadenopathy
    \item No pleural or pericardial effusion
    \item No distant metastases on the scanned field
\end{itemize}

\subsection{Differential Diagnosis}
\begin{table}[h]
\centering
\caption{Differential for hilar mass in a smoker}
\begin{tabular}{|p{0.38\textwidth}|p{0.50\textwidth}|}
\hline
\textbf{Diagnosis} & \textbf{Key features} \\
\hline
Hilar lymphoma & Bilateral adenopathy, younger patients \\
\hline
Hilar tuberculosis & Calcified nodes, upper lobe cavitary disease \\
\hline
Lung carcinoma & Spiculated mass, lymphadenopathy, smoking history \\
\hline
Post-obstructive pneumonia & Collapse distal to endobronchial lesion \\
\hline
Sarcoidosis & Bilateral hilar lymphadenopathy, non-smokers \\
\hline
\end{tabular}
\end{table}

\subsection{Final Diagnosis}
Non-small cell lung cancer (squamous cell carcinoma), T2aN2M0---Stage IIIA.

\subsection{Management}
PET-CT for staging, followed by neoadjuvant chemoradiotherapy and surgical review.

\subsection{Teaching Point}
A lung lesion that fails to resolve after 6 weeks of appropriate antibiotic therapy must be investigated as a potential malignancy. CT with contrast is the next step. Key CT features favouring malignancy include spiculation, size $>3$ cm, irregular margins, and associated lymphadenopathy. Always compare with prior imaging to assess growth rate.

\clearpage

% ========================================================================
\section{Case 3: ACL Tear on Knee MRI}
% ========================================================================

\subsection{Clinical Presentation}
A 22-year-old football player presents with a swollen, painful left knee after a non-contact pivoting injury during a match. He felt a ``pop'' and was unable to continue playing. Lachman test is positive with a soft endpoint.

\subsection{Investigation}
MRI of the left knee without contrast (3T, dedicated knee coil).

\subsection{Imaging Findings}
\begin{itemize}
    \item Complete disruption of the anterior cruciate ligament fibres with high signal throughout the ligament
    \item The ACL is not taut and shows a wavy, irregular contour
    \item Moderate joint effusion
    \item Bone bruise pattern on the lateral femoral condyle and posterolateral tibial plateau (characteristic of pivot-shift mechanism)
    \item Intact posterior cruciate ligament, medial and lateral collateral ligaments
    \item No bucket-handle meniscal tear
\end{itemize}

\subsection{Differential Diagnosis}
\begin{table}[h]
\centering
\caption{Differential for acute knee injury with effusion}
\begin{tabular}{|p{0.34\textwidth}|p{0.54\textwidth}|}
\hline
\textbf{Diagnosis} & \textbf{MRI features} \\
\hline
ACL tear & Disrupted fibres, bone bruise pattern, positive Lachman \\
\hline
MCL sprain & Medial collateral ligament thickening, high signal \\
\hline
Meniscal tear & Linear high signal extending to articular surface \\
\hline
Osteochondral fracture & Defect in articular cartilage and subchondral bone \\
\hline
PCL tear & Disruption of posterior cruciate ligament \\
\hline
\end{tabular}
\end{table}

\subsection{Final Diagnosis}
Complete rupture of the anterior cruciate ligament.

\subsection{Management}
Referral to orthopaedic surgery for ACL reconstruction. Physiotherapy and rehabilitation plan initiated.

\subsection{Teaching Point}
The characteristic bone bruise pattern (lateral femoral condyle and posterolateral tibial plateau) reflects the pivot-shift mechanism and is a secondary sign of ACL rupture. On MRI, a normal ACL appears as a taut, low-signal band oriented at approximately 45$^\circ$ to the femoral shaft. Complete tears show discontinuity and high signal replacing the normal ligament.

\clearpage

% ========================================================================
\section{Case 4: Ovarian Torsion on Ultrasound}
% ========================================================================

\subsection{Clinical Presentation}
A 19-year-old woman presents with sudden-onset, severe left-sided pelvic pain with vomiting. She has a history of a known 6 cm left ovarian cyst on ultrasound 3 months ago. Pain is constant and worse with movement.

\subsection{Investigation}
Transvaginal and transabdominal ultrasound with color Doppler.

\subsection{Imaging Findings}
\begin{itemize}
    \item Enlarged left ovary (8 $\times$ 6 $\times$ 5 cm) with peripheral follicles
    \item The 6 cm simple cyst is now eccentrically located within the ovary
    \item Ovarian stroma is swollen and echogenic (oedematous)
    \item Color Doppler shows absent arterial flow within the ovary
    \item Small amount of free fluid in the pouch of Douglas
    \item Right ovary appears normal with normal Doppler flow
\end{itemize}

\subsection{Differential Diagnosis}
\begin{table}[h]
\centering
\caption{Differential for acute pelvic pain in a young woman}
\begin{tabular}{|p{0.39\textwidth}|p{0.49\textwidth}|}
\hline
\textbf{Diagnosis} & \textbf{Key ultrasound features} \\
\hline
Appendicitis & RLQ mass, non-compressible appendix \\
\hline
Ectopic pregnancy & Adnexal mass, free fluid, positive $\beta$hCG \\
\hline
Ovarian torsion & Enlarged ovary, absent Doppler flow, free fluid \\
\hline
Pelvic inflammatory disease & Thickened fallopian tubes, tubo-ovarian abscess \\
\hline
Ruptured ovarian cyst & Complex free fluid, collapsed cyst \\
\hline
\end{tabular}
\end{table}

\subsection{Final Diagnosis}
Left ovarian torsion.

\subsection{Management}
Emergency gynaecological consultation. Diagnostic laparoscopy performed within 2 hours. Ovary found to be ischemic but viable; cystectomy performed with ovarian detorsion.

\subsection{Teaching Point}
Absence of Doppler flow is highly specific but not always present in torsion (partial torsion may preserve some flow). The key ultrasound finding is an enlarged, oedematous ovary with peripheral follicles (``string of pearls'' sign). The clinical scenario---sudden onset, severe pain, known cyst---is as important as the imaging. Do not delay surgical consultation while awaiting imaging if clinical suspicion is high.

\clearpage

% ========================================================================
\section{Case 5: Liver Hemangioma on MRI}
% ========================================================================

\subsection{Clinical Presentation}
A 45-year-old woman is found to have a liver lesion on ultrasound performed for vague right upper quadrant discomfort. Liver function tests are normal. No history of malignancy. Tumor markers (AFP, CEA, CA19-9) are within normal limits.

\subsection{Investigation}
MRI liver with extracellular IV gadolinium contrast (or multiphasic contrast-enhanced MRI).

\subsection{Imaging Findings}
\begin{itemize}
    \item 3 cm well-circumscribed lesion in segment VI of the liver
    \item T2-weighted images: uniformly very bright (``light bulb'' sign)
    \item Pre-contrast T1: low signal
    \item Arterial phase: peripheral nodular discontinuous enhancement (``peripheral puddling'')
    \item Portal venous phase: progressive centripetal fill-in
    \item Delayed phase (extracellular contrast, 3--5 min): persistent homogeneous or near-complete fill-in (note: if a hepatocyte-specific agent is used, haemangiomas lack functioning hepatocytes and appear \emph{hypointense} on the 20-minute hepatocyte phase, unlike focal nodular hyperplasia)
\end{itemize}

\subsection{Differential Diagnosis}
\begin{table}[h]
\centering
\caption{Differential for a T2-bright liver lesion}
\begin{tabular}{|p{0.32\textwidth}|p{0.56\textwidth}|}
\hline
\textbf{Diagnosis} & \textbf{MRI features} \\
\hline
Focal nodular hyperplasia & Central scar, arterial hyperenhancement, isointense on hepatocyte phase \\
\hline
Hemangioma & T2-bright, peripheral nodular enhancement, centripetal fill \\
\hline
Hepatic adenoma & Homogeneous enhancement, associated with OCP use \\
\hline
Hepatocellular carcinoma & Washout, capsule, risk factors (cirrhosis, hepatitis) \\
\hline
Hypervascular metastasis & Rim enhancement, washout, known primary \\
\hline
\end{tabular}
\end{table}

\subsection{Final Diagnosis}
Hepatic hemangioma---no treatment required.

\subsection{Management}
Reassurance and routine follow-up. No biopsy needed given characteristic imaging features. Patient counseled that this is a benign vascular lesion.

\subsection{Teaching Point}
The combination of marked T2 hyperintensity (light bulb sign) and peripheral nodular discontinuous enhancement with progressive centripetal fill-in on delayed extracellular-contrast phases is essentially pathognomonic for hemangioma. No biopsy is needed. This is the most common benign liver tumor, found in up to 20\% of the population, and requires no treatment unless symptomatic.

\clearpage

% ========================================================================
\section{Case 6: Occult Hip Fracture on X-ray and MRI}
% ========================================================================

\subsection{Clinical Presentation}
An 81-year-old woman presents with right hip pain after a low-energy fall at home. She can bear partial weight but has significant pain on internal rotation. Initial X-rays are reported as normal.

\subsection{Investigation}
Initial: AP pelvis and lateral hip X-ray. Follow-up: MRI right hip without contrast.

\subsection{Imaging Findings}
\textbf{X-ray (initial report: normal):}
\begin{itemize}
    \item No fracture line visible
    \item No cortical disruption or step-off
    \item No joint effusion visible
    \item Mild background osteopenia
\end{itemize}

\textbf{MRI (ordered due to persistent clinical suspicion):}
\begin{itemize}
    \item Linear low signal through the femoral neck on T1-weighted images (fracture line)
    \item Corresponding high signal on STIR sequences (bone marrow edema)
    \item No displacement
    \item Small joint effusion
    \item Intact femoral head vascularity (no avascular necrosis)
\end{itemize}

\subsection{Differential Diagnosis}
\begin{table}[h]
\centering
\caption{Differential for hip pain after a fall in elderly patients}
\begin{tabular}{|p{0.42\textwidth}|p{0.46\textwidth}|}
\hline
\textbf{Diagnosis} & \textbf{Imaging features} \\
\hline
Hip dislocation & Joint incongruity, obvious on X-ray \\
\hline
Intertrochanteric fracture & Visible fracture on X-ray \\
\hline
Occult femoral neck fracture & Normal X-ray, MRI-positive fracture line \\
\hline
Pelvic fracture & Fracture through pelvic ring \\
\hline
Pre-existing OA flare & Joint space narrowing, osteophytes \\
\hline
\end{tabular}
\end{table}

\subsection{Final Diagnosis}
Non-displaced occult fracture of the right femoral neck.

\subsection{Management}
Orthopaedic consultation. Internal fixation with cannulated screws to prevent displacement. Non-weight-bearing precautions for 6 weeks.

\subsection{Teaching Point}
Up to 2--3\% of elderly patients with hip pain after a fall and normal X-rays will have an occult fracture on MRI. If clinical suspicion remains high despite normal radiographs, MRI is the investigation of choice (sensitivity $>99\%$). CT has lower sensitivity for non-displaced femoral neck fractures. Never dismiss persistent hip pain in an elderly patient based on a normal X-ray alone.

\clearpage
